\documentclass[11pt,a4paper]{article}

% ------------------------------------------------------------------
% House layout, shared with the sibling preprints in docs/research/:
% 11pt a4, 2.5cm margins, Computer Modern, \maketitle front matter,
% abstract + keywords + JEL codes, numbered sections, table of contents.
% ------------------------------------------------------------------
\usepackage{amsmath, amssymb, amsthm}
\usepackage[a4paper, margin=2.5cm]{geometry}
\usepackage{array}
\usepackage{longtable}

\newcommand{\R}{\mathbb{R}}
\newcommand{\E}{\mathbb{E}}
\DeclareMathOperator{\IC}{IC}
\newcommand{\code}[1]{\texttt{\detokenize{#1}}}

% booktabs unavailable in this distribution — plain equivalents
\newcommand{\toprule}{\hline\hline}
\newcommand{\midrule}{\hline}
\newcommand{\bottomrule}{\hline\hline}
\usepackage[round]{natbib}

\begin{document}

% ============================================================

\title{\textbf{Fifteen Mechanisms That Did Not Work\\[3pt]
on a Perpetual-Futures DEX}\\[6pt]
\large A Negative-Results Census, with Reproducible
Microstructure Measurements}

\author{Yigit Onat\thanks{Working paper, preliminary draft --- comments
welcome. All estimates are from an internal empirical record
consolidated 2026-08-12; section references of the form \S$n$ in tables
refer to that record and are retained for provenance. Please do not
cite specific numbers without checking for a newer version.}\\
\textit{Independent Researcher}\\
\texttt{yionat@gmail.com}}
\date{August 2026}
\maketitle

% ============================================================
\begin{abstract}
Over six months and roughly ten thousand tested configurations on
Hyperliquid, a fully on-chain perpetual-futures exchange, we attempted
to monetise short-horizon predictive structure in the limit order book
and adjacent mechanisms, and failed fifteen times. This paper reports
the graveyard deliberately: negative results in trading research are
rarely published, so each generation of researchers re-derives the same
walls at full cost. We organise the record by \emph{economic mechanism}
(who is on the other side, and why they would lose) and by \emph{death
reason} (arithmetic, mechanism absent, underpowered, or blocked), a
distinction that decides whether a question is closed or merely
unresolved. The common signature of every refutation is stated and
quantified: the predicted move is comparable to the round-trip cost.
Alongside the graveyard we report descriptive measurements of the venue
that are, to our knowledge, undocumented in the academic literature:
the median pair quotes a half-spread $17$--$19\times$ that of BTC;
spread and depth are uncorrelated in cross-section ($\rho=-0.107$); the
median pair carries under \$100 of notional at the touch and only 4 of
177 pairs hold \$5{,}000; funding settles hourly; and the venue's
candle history is capped at ${\sim}5000$ bars per interval, making
fine-grained history unrecoverable unless accumulated live. A frozen
data slice and a one-command script reproducing the headline
measurements ship with the paper. We close with the transferable
result: a twenty-minute expected-move-to-cost calculation, run before
any infrastructure is built, would have redirected the entire programme
in week one.
\end{abstract}

\smallskip
\noindent\textbf{Keywords:} negative results, market microstructure,
perpetual futures, decentralised exchanges, transaction costs, adverse
selection, reproducibility

\smallskip
\noindent\textbf{JEL Classification:} C58, G12, G14, G23

\tableofcontents

% ============================================================
\section{Introduction}
\label{sec:intro}

The empirical starting point of this programme was a strong positive
result, documented in a companion paper: order-book imbalance on BTC,
ETH and SOL perpetuals predicts the direction of the mid-price with a
rank information coefficient of ${\approx}0.45$ at horizons of one to
five seconds, universally across symbols and volatility regimes. The
result survives every statistical control we ran at it. What follows
here is the record of attempting to turn that structure, and fourteen
adjacent mechanisms, into realised profit --- and failing every time.

Negative results of this kind are close to unpublishable in the
standard venues, which is precisely why they are worth writing down.
The walls documented below are structural properties of the venue and
of transaction-cost arithmetic; they do not depend on our estimators,
and anyone entering this market will hit them at full cost unless told
where they are. Fifteen refutations at roughly ten thousand tested
configurations is, in effect, a large prepaid experiment whose output
is a map.

Three choices shape the paper:

\begin{enumerate}
\item \textbf{Organisation by mechanism, not by technique.} Hawkes
intensities, order-flow imbalance, VPIN, jump detection, OU
mean-reversion and book-pressure features are six different techniques
--- and all six draw on one economic source: short-horizon order-flow
pressure in a single venue's book. Organised by technique, the record
looks diversified; organised by \emph{who pays}, the concentration is
visible, and the failures stop being surprising. Two techniques applied
to one edge source are not diversification.
\item \textbf{Death reasons are first-class.} ``This failed'' is not
one statement but four (Section~\ref{sec:deaths}), and conflating them
is how live questions get buried and dead ones exhumed. Each entry in
the census carries its death reason, and only two of the four reasons
close a question permanently.
\item \textbf{The measurements are the contribution; the graveyard is
the service.} The descriptive facts in Section~\ref{sec:venue} ---
spread structure, depth, funding cadence, history retention, and the
survivorship boundary of the data itself --- hold independently of any
trading hypothesis, are reproducible from the frozen slice shipped with
this paper (Appendix~\ref{app:repro}), and several of them contradict
intuitions imported from equity microstructure or from centralised
crypto venues.
\end{enumerate}

One boundary is deliberate: this paper publishes methods, measurements
and refutations. It does not publish anything the record currently
treats as a live, unrefuted edge, and the census below should be read
with that filter in mind.

\section{The venue and the data}
\label{sec:data}

Hyperliquid is a fully on-chain perpetual-futures exchange with a
central limit order book, permissionless API access over REST and
WebSocket, an hourly funding clock
(Section~\ref{sec:venue:funding}), and --- unusually ---
publicly observable wallet positions and liquidation levels. At the
time of the census the venue listed 177 tradeable perpetual pairs.

The internal platform ingests the venue's book and trade stream at
100\,ms cadence into a 236-feature representation per symbol (BTC, ETH,
SOL; ${\sim}2.17$ million ticks per symbol in the core sample), and
maintains three slower archives across the full universe: OHLCV candles
at 1m/5m/15m/1h (${\sim}3.06$ million candles, 708 series), periodic
full-book snapshots for all 177 pairs, and, since August 2026, the
hourly funding-rate history of every pair.

\paragraph{Survivorship, declared.} The candle archive covers the 177
pairs listed during the collection window. At least 55 previously
listed perpetuals were delisted before or during the window and are
absent. Delisting concentrates in exactly the thin, low-attention tail
where several of the untested mechanisms (notably gradual information
diffusion, Section~\ref{sec:census}) predict their effect, so
cross-sectional results below are conditioned on survival to an extent
we cannot fully quantify. We state this rather than let a referee find
it: every cross-sectional number in this paper should be read as
``among pairs that survived to be measured.''

\paragraph{History is not recoverable.} The venue caps retrievable
candle history at ${\sim}5000$ bars per interval regardless of start
time. One-minute history therefore expires ${\sim}3.5$ days after the
fact and \emph{cannot be backfilled at any later date}; the same
applies pro-rata to every interval. There is, to our knowledge, no
public archive of this venue's L2 book states. Datasets here are
accumulated live or lost, which is why the reproducibility slice
(Appendix~\ref{app:repro}) is frozen into the repository rather than
re-fetched on demand.

\section{Descriptive measurements}
\label{sec:venue}

These facts constrain every strategy on the venue, are independent of
any signal, and each one killed or reshaped at least one hypothesis in
the census.

\subsection{The universe is not its majors}
\label{sec:venue:spread}

Measured over repeated full-book snapshot sweeps of all 177 pairs, the
cross-sectional median half-spread is $1.37$--$1.43$ basis points
against BTC's $0.077$--$0.083$\,bps: \textbf{the median pair quotes
$17$--$19\times$ wider than BTC}, and 169 of 177 pairs are wider than
the three majors on which all high-frequency work was done. Any result
established on BTC's book generalises to almost nowhere on the venue,
and any cost model calibrated on the majors understates the universe by
an order of magnitude.

\subsection{Depth does not follow spread}
\label{sec:venue:depth}

Across the same sweeps, per-pair half-spread and touch depth are
essentially uncorrelated in cross-section ($\rho = -0.107$): wide pairs
are not deep pairs resting at a fairer price, they are simply thin. The
median pair carries \$33--\$391 of notional at the best quotes
(slice median \$81), and \textbf{only 4 of 177 pairs hold \$5{,}000 at
the touch}. The extreme cases are instructive: pairs quoting
$13$--$27$\,bps half-spreads over \$3--\$20 of touch notional. A large
per-fill edge on \$20 of size is not a business. Admission to any
cross-sectional strategy is therefore a \emph{joint} spread--depth
condition, and the joint condition is far more restrictive than either
margin: a 2\,bps half-spread ceiling alone admits 129 pairs, but the
capacity to deploy even \$1{,}000 per pair per day at 1\% of volume
cuts the tradeable set to ${\sim}117$, and \$5{,}000 of touch size
exists in four names.

\subsection{Funding settles hourly, and the displayed rate is hourly}
\label{sec:venue:funding}

Perpetual funding on this venue is exchanged \textbf{every hour}
(verified directly: settlement timestamps in the venue's own funding
history are spaced $1.0000$\,h apart), and the funding field in the
live asset context is the hourly rate --- observed baseline
$1.25\times10^{-5}$ per hour, i.e.\ the venue's $0.00125\%$/h default.
Conventional CEX intuition (8-hour funding) mis-prices held positions
on this venue by a factor of eight. In our own simulation stack this
exact error survived undetected for months --- funding was charged in
\emph{no} backtest, and the configured interval was 8 hours --- an
error we document in Section~\ref{sec:method} because the manner of
finding it generalises.

\subsection{Retention, cadence, and what they imply}

The ${\sim}5000$-bar retention cap (Section~\ref{sec:data}) has a
strategic consequence that is easy to miss: for any research programme
on this venue, \emph{data continuity is the binding constraint}, ahead
of ideas, compute, or capital. A gap in live collection is permanent.
Half of the operational machinery described in
Section~\ref{sec:method} exists to protect continuity, and the single
most damaging defect we record (Section~\ref{sec:census},
entry~M) was an availability failure, not a statistical one.

\section{Method, and what the method caught}
\label{sec:method}

The statistical equipment is standard; what we would emphasise is the
combination, because in every near-miss of the final month it was a
\emph{control}, not judgement, that killed the hypothesis.

\begin{itemize}
\item \textbf{Planted-signal tests before real data.} Every estimator
must first recover a synthetic effect of known size planted in
synthetic data, and return null on planted noise. Three separate
estimator bugs --- each of which would have manufactured a discovery
--- were caught only this way.
\item \textbf{Null calibration under the right exchangeability.}
Permutation nulls are constructed to preserve the dependence structure
the statistic is exposed to (day-block permutations for serial
dependence, label permutations that preserve group sizes for
cross-sectional claims). One internal audit produced an IC of
$0.39$--$0.46$ --- \emph{higher} than the claim it was auditing ---
purely from overlapping forward windows; non-overlapping sampling is
mandatory since.
\item \textbf{False-discovery control across the whole grid.}
Benjamini--Hochberg at $q=0.05$ over every cell actually computed, not
over the cells that looked interesting.
\item \textbf{Per-day durability, not pooled means.} A pooled mean
over 25 days is routinely carried by two days. Every claim must hold on
a stated fraction of individual days.
\item \textbf{Pre-registration in version control.} Acceptance
criteria, trial counts for deflated-Sharpe penalties
\citep[in the sense of][]{bailey2014}, and stop rules are committed to
the repository \emph{before} a study runs; the commit hash is the
proof. The funding-carry study of entry~O below is the worked example:
its gate, criteria, and sceptical priors were committed before its
data was fetched, and the study's own gate was not permitted to be
adjusted afterwards.
\item \textbf{Cost realism as a gate, not a sensitivity.} All costs
flow from one configuration file through one loader. The census
contains a five-fold false-discovery event (entry~C) whose root cause
was a silently wrong venue's fee tier reaching a harness default.
\end{itemize}

Two of the fifteen entries below are methodological deaths --- the
hypothesis was never really tested because the harness was wrong --- and
we count them in the census deliberately. A negative-results paper that
omitted its own instrumentation failures would be reporting a cleaner
record than the one that exists.

\section{How a mechanism dies}
\label{sec:deaths}

\begin{table}[h]
\centering
\begin{tabular}{p{4.2cm} p{2.2cm} p{7.4cm}}
\toprule
Death reason & Permanent? & Implication \\
\midrule
\textbf{Arithmetic} --- predicted move $\approx$ cost & Yes, at that
horizon & Only a cost change reopens it \\
\textbf{Mechanism absent} --- measured, not there & Yes & Closed \\
\textbf{Undecidable} --- underpowered & No & Revisit at adequate $n$; a
\emph{collection} task, not a research one \\
\textbf{Blocked} --- data or access missing & No & Acquire the data;
the question is untested, not answered \\
\bottomrule
\end{tabular}
\caption{The four death reasons. Only the first two close a question.
Filing an underpowered result as ``refuted'' discards a live question;
filing an arithmetic death as ``needs more data'' wastes a year. The
census tags every entry with one of these four.}
\label{tab:deaths}
\end{table}

The distinction earns its keep twice in this record: two of the
strongest-looking refutations (entries~J and~L) are formally
\emph{undecidable} --- one observed $1.5$ events per day against a
design that needed hundreds, the other needs ${\sim}0.9$ years of
rebalances for its Sharpe ratio to reach significance --- and treating
either as ``mechanism absent'' would have been a classification error
with multi-month consequences in both directions.

\section{The census}
\label{sec:census}

Table~\ref{tab:census} is the record. Family numbers refer to the
economic taxonomy of Section~\ref{sec:families}.

{\small
\begin{longtable}{p{0.5cm} p{5.2cm} p{3.8cm} p{1.4cm} p{2.9cm}}
\caption{The fifteen-entry census. ``Family'' refers to
Table~\ref{tab:families}. Entries C and M are instrumentation deaths
and are included deliberately (Section~\ref{sec:method}).}
\label{tab:census}\\
\toprule
\# & What was tried & Outcome & Family & Death reason \\
\midrule
\endfirsthead
\toprule
\# & What was tried & Outcome & Family & Death reason \\
\midrule
\endhead
A & Taker capture of the 1--5\,s imbalance signal & $0.5$--$2$\,bps
predicted move vs $11$\,bps round trip & 1/11 & Arithmetic \\
B & Naive maker capture of the same signal & $\IC$ $0.45 \to 0.03$
conditional on directionally-consistent fills & 1 & Arithmetic (adverse
selection) \\
C & Five production signal ``winners'' & All five refuted on re-audit
& --- & Methodological (wrong venue's fee tier in harness defaults) \\
D & VWAP-reversion at taker cost, 58 days & $-25$k to $-41$k bps
cumulative & 1/11 & Arithmetic (${\sim}50$ trades/day $\times$ fees) \\
E & Textbook Avellaneda--Stoikov quoting & $-1.89$\,bps/fill;
risk-widened quotes rest behind the touch and fill only through
adverse moves & 1 & Model-transfer failure \\
F & Touch-pegged passive quoting, 8 configs $\times$ 179
symbol-days & 0 survive day-consistency and concentration & 1 &
Mechanism absent (at the majors' touch) \\
G & Fee-tier staking discounts as the fix & Verdict tier-invariant:
discounts apply to fees, never rebates & 1 & Arithmetic \\
H & The full maker rebate ladder & Breakeven rebate $+0.144$\,bps
exceeds the best volume-tier rate (zero); first viable rung requires
${\sim}1.5\%$ of venue-wide maker volume & 1 & Arithmetic
(structural) \\
I & Momentum persistence, 1m/5m bars, majors & Anti-persistent: 34/36
cells negative; the universe mean-reverts ${\sim}$8-to-1 & 9 &
Mechanism absent (naked signature, bar scale) \\
J & VWAP-band oscillation harvesting, 330 cells & 0 survive --- but
${\sim}1.5$ events/day against a design needing hundreds & 1 &
\emph{Undecidable} \\
K & Permutation entropy as cross-sectional score & Interquartile range
$0.0005$: the middle half of the universe is indistinguishable & --- &
Mechanism absent \\
L & Cross-sectional low-vol rotation (long/short, out of sample) & 0/6
pre-registered criteria; beta-neutral rebuild reaches 4/6 with Sharpe
2.1 but needs ${\sim}0.89$\,yr for significance & 8/9 &
\emph{Undecidable} \\
M & Hierarchical signal combiner & Loses to one of its own inputs;
weights were fitted after the window they were scored on & --- &
Methodological \\
N & Agreement gating (trade only when signals agree) & Harmful: the
\emph{disagreement} subset carries the signal & 2 & Mechanism absent
(as specified) \\
O & Funding carry, dollar-neutral across 118 pairs, 92 days,
pre-registered & Funding leg collects $+5$--$11\%$ gross; the
short-the-crowd price leg loses more; 0/6 configs pass & 4 & Mechanism
absent (carry $\approx$ fair drift compensation) \\
\bottomrule
\end{longtable}}

Three entries deserve prose beyond the table.

\paragraph{B: the execution boundary.} The strongest structure in the
record --- directional $\IC \approx 0.45$ at seconds --- collapses to
${\approx}0.03$ when conditioned on obtaining a directionally-correct
passive fill, while conditioning on the presence of \emph{any} fill
\emph{raises} the coefficient to $0.52$. The information is real and
the market microstructure prices it: a passive order fills
preferentially when the move is against it. This isolates adverse
selection, not signal decay, as the binding constraint, and it is the
single result that most reshaped the programme.

\paragraph{H: zero fees are not free money.} At BTC's touch the
half-spread ($0.077$--$0.083$\,bps) is roughly one third of the
expected adverse move conditional on a fill ($0.23$--$0.24$\,bps), so a
resting quote loses ${\sim}0.08$\,bps per posting even at a zero fee
rate. The venue must \emph{pay} ${\sim}+0.15$\,bps before passive
quoting at the majors' touch breaks even, and the first rebate rung
that clears this sits at ${\sim}1.5\%$ of venue-wide 14-day maker
volume (${\sim}\$86$M/day of maker fills at current venue volume) ---
an entry ticket, not a fee schedule. The reproduction script regenerates
this ladder from the frozen slice and the venue's published tiers.

\paragraph{O: the carry that was exactly fair.} The funding-carry study
is the census's worked example of the full method: mechanism named in
advance (the crowded side pays hourly; a dollar-neutral book collects
the spread), gate and criteria committed to version control before the
data existed, and a cost-ratio gate (expected carry over expected
churn cost $\ge 3$) evaluated \emph{before} the backtest was permitted
to run. The gate passed at $4.5$; the book still failed. Funding was
collected almost exactly as forecast ($+5$--$11\%$ gross over 92 days)
and the price leg gave it back with interest: the coins whose crowd
paid the most kept moving the crowd's way. On this venue, over this
window, funding was a fair price for drift --- a textbook risk-premium
story, measured rather than assumed. The two marginally positive slow
configurations are disqualified rather than suggestive: the best
carries 92\% of its P\&L in one day and would need
${\sim}4.6$ years for its Sharpe ratio to clear significance.

\section{The mechanism families, and what remains open}
\label{sec:families}

{\small
\begin{longtable}{p{0.5cm} p{5.6cm} p{7.6cm}}
\caption{Families by economic source. Numbers are stable identifiers
in the internal taxonomy; 6, 10, 12 were re-specified or retired and
never reused --- gaps are deliberate. Two families are omitted as
operationally inaccessible (volatility risk premium: no options
market; queue-priority/latency: no colocation).}
\label{tab:families}\\
\toprule
\# & Family --- economic source & Status after the census \\
\midrule
\endfirsthead
\toprule
\# & Family --- economic source & Status after the census \\
\midrule
\endhead
1 & Liquidity provision --- paid to bear inventory + adverse selection
\citep[cf.][]{grossman1988, avellaneda2008} & Heavily worked; closed at
the majors' touch by arithmetic (entries E--H); open on wide pairs,
where the half-spread is a multiple of BTC's \\
2 & Adverse-selection avoidance \citep[cf.][]{kyle1985, easley2012} &
Survives as a \emph{gate} only: flow-toxicity conditioning lifts
per-fill economics on all three majors but carries no direction \\
3 & Price-discovery lag across venues \citep[cf.][]{hasbrouck1995,
makarov2020} & Untried --- blocked on a cross-venue feed \\
4 & Funding reflexivity & Carry branch measured and dead (entry O);
the reflexive-unwinding branch (funding $\to$ forced flow $\to$ price)
untested, blocked on position-history accrual \\
5 & Deterministic liquidation --- the engine executes at a computable
price \citep[nearest analogue][]{coval2007} & Best-supported untried
family: trigger prices are \emph{public} on this venue; blocked on
instrumentation, not theory \\
6 & Gradual information diffusion \citep[cf.][]{hong1999, hong2000} &
Untried at its native horizon (weeks); the minute-scale reversion
result does not touch it, and the pooled tests that exist would dilute
a liquidity-conditional effect to zero by construction \\
7 & Attention and flow-driven demand \citep[cf.][]{barber2008} &
Untried; listing events exist in-archive but $n=2$ observable events
so far --- a collection task \\
8 & Statistical relative value \citep[cf.][]{gatev2006} & Registered,
never evaluated \\
9 & Behavioural under/over-reaction \citep[cf.][]{jegadeesh1993,
debondt1985} & The \emph{naked signature} is exhausted at bar scale
(entry I); the generating mechanisms are largely untried \\
11 & Microstructure noise / bid--ask bounce \citep[cf.][]{roll1984} &
Exhausted \\
\bottomrule
\end{longtable}}

Two remarks the table cannot carry.

\paragraph{Signatures are not mechanisms.} Momentum and mean-reversion
are signs of autocorrelation --- observables, not explanations. At
least four distinct mechanisms produce reversion (inventory
compensation, forced-flow overshoot, bid-ask bounce, behavioural
over-reaction) and four produce momentum (discovery lag, gradual
diffusion, attention flow, under-reaction). What this record exhausted
is the \emph{naked signature} --- trade the sign of recent returns, at
bar scale, at taker cost --- which says nothing about the generating
mechanisms individually. The question to ask any momentum or reversion
proposal is not ``what horizon?'' but ``\emph{paid by whom?}''

\paragraph{The strongest negative is a signpost.} ``The universe
reverts 8-to-1'' (entry I) is precisely the condition under which
liquidity provision and forced-flow harvesting pay, because both are
compensated \emph{by} reversion. The most robust refutation in the
record therefore points directly at the two families the venue's
architecture makes most observable --- and both are blocked on
instrumentation rather than closed. On-chain position visibility is a
structural research advantage of this venue class: forced flow that
equity-market studies must infer from returns is, here, directly
observable in advance.

\section{The transferable result}
\label{sec:lesson}

Every arithmetic death in Table~\ref{tab:census} was knowable in
advance. The predicted move of the core signal ($0.5$--$2$\,bps at
seconds) against the realistic round trip ($\sim$$11$\,bps taker) is a
ratio of ${\sim}0.1$; the programme's subsequently adopted rule ---
\emph{expected move over round-trip cost} $\ge 3$ \emph{at the design
horizon, computed before anything is built} --- would have failed the
entire high-frequency taker programme on day one, using twenty minutes
and public data. The same twenty minutes, run at gate zero of the
funding-carry study (entry O), correctly admitted it --- and the study
then failed for a reason the ratio cannot see, which is the right
division of labour: the ratio filters what cannot work; measurement
decides what does not.

Stated for the general case: on any venue, for any signal, compute
$m/c$ where $m$ is the expected favourable move over the holding
horizon and $c$ the full round trip at the fee tier you actually hold.
Below $1$, the strategy is arithmetically dead regardless of the
information content; between $1$ and $3$, it is a fill-quality
research programme, not a trading strategy; above $3$, it earns the
right to a backtest. The census contains six entries that this
one-line calculation would have prevented.

\section*{Acknowledgements}
The platform, studies and record summarised here were built and run by
the author. No external funding; no conflicts to declare.

% ============================================================
\appendix

\section{Reproducibility}
\label{app:repro}

The repository ships a frozen public slice and a one-command
reproduction path:

\begin{center}
\code{./reproduce.sh}
\end{center}

\noindent requiring only Python with \code{pandas}, \code{pyarrow} and
\code{matplotlib}; no network access and no credentials. The script
reads the frozen slice (two full-universe L2 snapshot days --- 177
pairs, ${\sim}50$k book states, ${\sim}9.9$\,MB --- plus hourly candles
for the majors) and the venue fee configuration, and regenerates:

\begin{enumerate}
\item the cross-sectional half-spread distribution
(Section~\ref{sec:venue:spread}): slice values --- median $1.428$\,bps,
$18.6\times$ BTC;
\item the touch-depth distribution (Section~\ref{sec:venue:depth}):
slice values --- median \$81; $4/177$ pairs at \$5{,}000;
\item the maker-ladder breakeven (entry H): slice values --- breakeven
rebate $+0.151$\,bps; first viable rung unchanged.
\end{enumerate}

\noindent together with a machine-readable \code{headlines.json}. The
slice's own values are pinned by unit tests, so any drift in the
computation is a build failure rather than a silent revision. Two
constants (expected adverse selection conditional on a fill at BTC's
touch, $0.228/0.242$\,bps bid/ask) enter the ladder figure from tick
data too large to freeze and are declared as frozen inputs in the
script's header. Because the venue caps history retention and serves no
historical order-book states (Section~\ref{sec:data}), freezing the
slice in the repository is the only mechanism by which these
measurements remain independently checkable.

\renewcommand{\refname}{References}
\begin{thebibliography}{99}\small

\bibitem[Avellaneda and Stoikov(2008)]{avellaneda2008}
Avellaneda, M.\ and S.\ Stoikov (2008). High-frequency trading in a
limit order book. \emph{Quantitative Finance} 8(3), 217--224.

\bibitem[Bailey and L\'opez de Prado(2014)]{bailey2014}
Bailey, D.\ and M.\ L\'opez de Prado (2014). The deflated Sharpe
ratio. \emph{Journal of Portfolio Management} 40(5), 94--107.

\bibitem[Barber and Odean(2008)]{barber2008}
Barber, B.\ and T.\ Odean (2008). All that glitters: the effect of
attention and news on the buying behavior of individual and
institutional investors. \emph{Review of Financial Studies} 21(2),
785--818.

\bibitem[Coval and Stafford(2007)]{coval2007}
Coval, J.\ and E.\ Stafford (2007). Asset fire sales (and purchases)
in equity markets. \emph{Journal of Financial Economics} 86(2),
479--512.

\bibitem[De Bondt and Thaler(1985)]{debondt1985}
De Bondt, W.\ and R.\ Thaler (1985). Does the stock market overreact?
\emph{Journal of Finance} 40(3), 793--805.

\bibitem[Easley et~al.(2012)]{easley2012}
Easley, D., M.\ L\'opez de Prado and M.\ O'Hara (2012). Flow toxicity
and liquidity in a high-frequency world. \emph{Review of Financial
Studies} 25(5), 1457--1493.

\bibitem[Gatev et~al.(2006)]{gatev2006}
Gatev, E., W.\ Goetzmann and K.\ Rouwenhorst (2006). Pairs trading:
performance of a relative-value arbitrage rule. \emph{Review of
Financial Studies} 19(3), 797--827.

\bibitem[Grossman and Miller(1988)]{grossman1988}
Grossman, S.\ and M.\ Miller (1988). Liquidity and market structure.
\emph{Journal of Finance} 43(3), 617--633.

\bibitem[Hasbrouck(1995)]{hasbrouck1995}
Hasbrouck, J.\ (1995). One security, many markets: determining the
contributions to price discovery. \emph{Journal of Finance} 50(4),
1175--1199.

\bibitem[Hong and Stein(1999)]{hong1999}
Hong, H.\ and J.\ Stein (1999). A unified theory of underreaction,
momentum trading, and overreaction in asset markets. \emph{Journal of
Finance} 54(6), 2143--2184.

\bibitem[Hong et~al.(2000)]{hong2000}
Hong, H., T.\ Lim and J.\ Stein (2000). Bad news travels slowly: size,
analyst coverage, and the profitability of momentum strategies.
\emph{Journal of Finance} 55(1), 265--295.

\bibitem[Jegadeesh and Titman(1993)]{jegadeesh1993}
Jegadeesh, N.\ and S.\ Titman (1993). Returns to buying winners and
selling losers: implications for stock market efficiency.
\emph{Journal of Finance} 48(1), 65--91.

\bibitem[Kyle(1985)]{kyle1985}
Kyle, A.\ (1985). Continuous auctions and insider trading.
\emph{Econometrica} 53(6), 1315--1335.

\bibitem[Makarov and Schoar(2020)]{makarov2020}
Makarov, I.\ and A.\ Schoar (2020). Trading and arbitrage in
cryptocurrency markets. \emph{Journal of Financial Economics} 135(2),
293--319.

\bibitem[Roll(1984)]{roll1984}
Roll, R.\ (1984). A simple implicit measure of the effective bid-ask
spread in an efficient market. \emph{Journal of Finance} 39(4),
1127--1139.

\end{thebibliography}

\end{document}
